\section{Introduction}

The decennial Census drives congressional redistricting, but the demographic changes
that will shape the 2030 map are already visible in current ACS and Census Bureau
projection data. Acting on this information now allows redistricting infrastructure
to be built, validated, and ready before the April 2031 Census data release.

This paper uses the Census Bureau's 2020 Vintage Population Projections --- the
most recent long-range projections available, extending through 2060 --- to project
2030 state populations and apply Huntington-Hill apportionment. The projections come
in three series (low, medium, high), reflecting different assumptions about fertility,
mortality, and net international migration.

The practical stakes: a state that gains seats (TX, FL) needs new bisection tree
designs; a state that loses seats (NY, CA) may see major district mergers that
challenge algorithmic compactness. The bisect pipeline's bisection tree structure
--- the ApportionRegions hierarchy --- is determined entirely by the seat count $k$
and its prime factorization. Changes in $k$ cascade through the entire tree design.

\paragraph{Data source.}
Census Bureau National Population Projections, 2020 Vintage: Table NP2020-T1
(state-level projections by year, all series). Available at
\texttt{census.gov/programs-surveys/popproj.html}. We extract the 2030 column
from the medium (middle) series for the main projections, and the low and high
series for confidence bounds.
